\documentclass[10pt,a4paper]{amsart}
\usepackage[margin=19mm]{geometry}
\usepackage{amsmath,amssymb,mathtools}
\usepackage[scr=rsfs,cal=euler]{mathalfa}
\usepackage{xfrac}
\usepackage[unicode,hidelinks,bookmarks=false]{hyperref}
\newcommand{\ie}{i.e.\ }
\newcommand{\cf}{cf.\ }
\newcommand{\resp}{respectively\ }
\newcommand{\Subsection}{\subsection}
\providecommand{\nobreakdash}{\nobreak\hbox{-}}
\setlength{\emergencystretch}{2em}
\multlinegap0pt
\usepackage{xcolor}
\usepackage{amssymb}
\usepackage{enumerate}
\usepackage{cancel}
\usepackage[normalem]{ulem}
\newtheorem{proposition}{Proposition}
\title{Pseudodifferential Operators and the Connes-Kasparov Isomorphism}
\author{Peter Debello and Nigel Higson}
\date{Source: arXiv:2502.14985v2 (CC BY 4.0)}
\begin{document}
\maketitle

\noindent\emph{Source and licence.} Pseudodifferential Operators and the Connes-Kasparov Isomorphism. Version arXiv:2502.14985v2 (CC BY 4.0), distributed by arXiv under the Creative Commons Attribution 4.0 International licence, http://creativecommons.org/licenses/by/4.0/, as recorded in the arXiv licence metadata for this exact version and shown on the abstract page https://arxiv.org/abs/2502.14985v2. Full licence notice: \texttt{licenses/arxiv-2502.14985-CC-BY-4.0.txt}.\par\smallskip

\noindent\textit{Editorial scope.} Counted unit: the article's proposition prop-l-2-versus-s-star-boundedness (source0.tex lines 836--843). The article's complete original proof of it is delivered with it (source0.tex lines 845--901, one \texttt{proof} environment of 57 lines). No mathematics is written, repaired or re-derived: the statement and the proof are the article's own lines, in the article's own notation, and no retained line is substituted or re-flowed.  No further source-local context is needed: every cross-reference of every retained line resolves inside the two delivered spans.  Nothing was omitted from any retained span, so the package carries no omission record.  No retained line contains a citation, so the wrapper carries no bibliography and no bibliography entry.  No retained line contains a figure environment, an image-inclusion command or drawing code, so no image is delivered and the package declares no asset dependency.  Editorial formatting only, in addition: an AMS \texttt{amsart} preamble that restates the article's own declarations and macros used by the retained lines, namely the environments \texttt{proposition} on separate counters, together with the standard packages those lines need.  The retained environments print in this document as Proposition 1 in order of appearance, whereas the article prints the same items with the numbering of its own full document.  The complete native source of the article as distributed in the arXiv e-print for this exact version is bundled unchanged as \texttt{sources/173616/source0.tex}.
\begin{proposition}
\label{prop-l-2-versus-s-star-boundedness}
  Every  properly supported, $G$-equivariant, order zero classical pseudodifferential operator 
\[
A \colon [C_c^\infty ( G)\otimes V_1]^K \longrightarrow  [C_c^\infty ( G)\otimes V_2]^K 
\]
extends to a bounded operator between $L^2$-spaces.
\end{proposition}

\begin{proof} 
Let $\varphi$ be a smooth and compactly supported function on $K \backslash G$ such that 
\[
\int_G  \varphi (xg) \, dg = 1 \quad \forall x\in K\backslash G,
\]
and let 
\[
A_e {=}  A\varphi  \colon [C_c^\infty ( G)\otimes V_1]^K \longrightarrow  [C_c^\infty ( G)\otimes V_2]^K .
\]
This is a compactly supported,  order zero classical pseudodifferential operator on $M{=}K\backslash G$, and therefore it is bounded in the $L^2$-operator norm. If we set $A_g = g(A_e)$, then 
\[
(A f)(x) = \int_G (A_g f)(x) \, dg \quad \forall x \in K \backslash G,\quad \forall f\in [C_c^\infty ( G)\otimes V_1]^K  .
\]
The operators $A_g$  are unitarily equivalent to one another, and therefore they have the same $L^2$-operator norm:
\[
\| A_g \| = \| A_e\| \quad  \forall g\in G.
\]
Let  $C\subseteq K \backslash G$ be a compact subset such that the  support of $A_e$ is included in 
$C{\times } C$, and let 
\[
S = \{ \, g\in G : Kg \in C\, \} .
\]
This is a compact subset of $G$, and if  $f\in[C_c^\infty ( G)\otimes V_1]^K $, $g\in G$ and $x\in K \backslash G$, then $(A_gf)(x) = 0$ unless $x\in Cg$. In other words, $(A_gf)(x) = 0$ unless $g\in x^{-1}S$. 
So by the Cauchy-Schwarz inequality,
\[
\begin{aligned}
\| A f\|_{L^2}^2 
    & = \int _M \Bigl \| \int_G (A_gf)(x)\, dg \Bigr\|_{V_2}^2 \, dx \\
    & \le \int _M \Bigl ( \int_{x^{-1} S} 1 \, dg\Bigr ) \Bigl (   \int_{x^{-1}S} \bigl \|(A_gf)(x)\bigr \|_{V_2} ^2 \, dg\Bigr )   \, dx \\
    & = \operatorname{vol}(S)\cdot   \int _M \Bigl (\int_{x^{-1}S} \bigl \|(A_gf)(x)\bigr \|_{V_2} ^2 \, dg \Bigr )   \, dx  .
\end{aligned}
\]
But now, using Fubini's theorem, we can write 
\[
\int _M \Bigl (\int_{x^{-1}S} \bigl \|(A_gf)(x)\bigr \|_{V_2} ^2 \, dg \Bigr )   \, dx \\
     = \int _G \Bigl (\int_{Cg}  \bigl \|(A_gf)(x)\bigr \|_{V_2} ^2 \, dx  \Bigr ) \, dg .
\]
The inner integral on the right may be estimated by
\[
\int_{Cg}  \bigl \|(A_gf)(x)\bigr \|_{V_2} ^2 \, dx \le \| A_e\|^2\cdot \int_{Cg}  \bigl \|f(x)\bigr \|_{V_1} ^2 \, dx ,
\]
and so,  by Fubini  once again,
\[
\begin{aligned}
\| A f\|_{L^2}^2  
    & \le  \operatorname{vol}(S)\cdot  \| A_e\|^2 \cdot 
    \int _G \Bigl (\int_{Cg}   \|f(x)  \|_{V_1} ^2 \, dx  \Bigr ) \, dg \\
    & = \operatorname{vol}(S)\cdot \| A_e\|^2  \cdot 
    \int _M \Bigl (\int_{x^{-1}S}   \|f(x)  \|_{V_1} ^2 \, dg \Bigr )   \, dx \\
    & = \operatorname{vol}(S)^2\cdot \| A_e\|^2  \cdot 
    \int _M    \|f(x)  \|_{V_1} ^2     \, dx
    \\
    & = \operatorname{vol}(S)^2 \cdot \| A_e\|^2 \cdot  \| f\|_{L^2}^2,
\end{aligned}
\]
which proves that $A$ is $L^2$-bounded, as required.
\end{proof} 

\end{document}
